%! Introduction to Polynomial Functions — Unit 2, Lesson 2.0 — Lesson Plan
\documentclass[10pt]{article}
\usepackage{precalculus-article}
\usepackage{precalculus-boxes}
\usepackage{graphicx}
\graphicspath{{images/}}
\newcommand{\TallMath}[1]{$\displaystyle #1\rule[-1.4em]{0pt}{3.2em}$}
\newcommand{\CourseName}{Precalculus}
\newcommand{\MeetingLength}{55 minutes}
\newcommand{\UnitNumberName}{Unit 2: Polynomial Functions \quad}
\newcommand{\LessonNumberName}{Lesson 2.0: Introduction to Polynomial Functions}

\begin{document}
\begin{center}
    {\Large\bfseries \CourseName \\
    {\small\bfseries \UnitNumberName \LessonNumberName}}
\end{center}

\begin{tcolorbox}[colback=lilac, colframe=plum, boxrule=0.9pt, arc=2mm,
    left=2mm, right=2mm, top=1.5mm, bottom=1.5mm]
    \textbf{Primary Objective:} Students will recognize a polynomial function from
    its written form, name its degree, leading term, leading coefficient, and
    constant term, write a polynomial in standard form, explain why expressions
    with negative or fractional exponents (or a variable in a denominator or
    exponent) are \textit{not} polynomials, and describe the family's graphs as
    smooth and unbroken.
    \hfill\textit{\small (Unit introduction --- no CED topic)}
\end{tcolorbox}

\renewcommand{\arraystretch}{1.6}
\begin{skillbox}[Priority Ideas \& Skills]{goldbox}
\begin{tabularx}{\linewidth}{@{}X X@{}}
\textbf{Priority Skills} &
\textbf{Key Understandings} \\
\midrule
Decide whether a given expression defines a polynomial function, and justify the
decision. &
\begin{itemize}[leftmargin=1.2em,itemsep=2pt,topsep=0pt]
  \item A polynomial is a sum of terms, each a real number times a whole-number
        power of $x$.
  \item Three disqualifiers: a negative or fractional exponent, a variable in a
        denominator, a variable in an exponent.
\end{itemize} \\
Identify the degree, leading term, leading coefficient, and constant term of a
polynomial, and write it in standard form. &
\begin{itemize}[leftmargin=1.2em,itemsep=2pt,topsep=0pt]
  \item Standard form orders terms by descending power, which puts the leading
        and constant terms at the two ends.
  \item The sign travels with the coefficient: $-x^3$ has leading coefficient
        $-1$.
\end{itemize} \\
Name the low-degree families and read turning points off a pre-drawn graph. &
\begin{itemize}[leftmargin=1.2em,itemsep=2pt,topsep=0pt]
  \item Degrees 0--4: constant, linear, quadratic, cubic, quartic.
  \item Polynomial graphs are smooth and unbroken; a degree-$n$ polynomial has at
        most $n-1$ turning points.
\end{itemize} \\
\end{tabularx}
\end{skillbox}

\begin{skillbox}[{Vocabulary, Concepts, \& Theorems}]{greenbox}
\renewcommand{\arraystretch}{1.4}
\begin{tabularx}{\linewidth}{@{} >{\leavevmode\bfseries\color{plum}}p{0.24\linewidth} X @{}}
Polynomial function & A sum of terms, each a real number times a whole-number power of $x$. \\
Degree              & The highest exponent on $x$ in the polynomial. \\
Leading term        & The term carrying the highest exponent; it comes first in standard form. \\
Leading coefficient & The number multiplying the leading term, minus sign included. \\
Constant term       & The term with no $x$; equivalently, the value at $x = 0$. \\
Standard form       & Terms written in descending order of exponent. \\
Turning point       & A point where the graph stops rising and starts falling, or the reverse. \\
\end{tabularx}
\end{skillbox}

\begin{skillbox}[Activate Prior Knowledge \& Spiral Review]{sky}
    Please see attached warm-up (spiral review).
    Skills reviewed: evaluating an algebraic expression containing exponents,
    combining like terms,
    evaluating a function at a given value using $f(x)$ notation,
    reading a value off a pre-drawn graph.
\end{skillbox}

\begin{skillbox}[Hook]{sky}
    \textbf{Which three belong together?} Put six expressions on the board and ask
    students to circle the three that ``belong to the same family'':
    \[
      3x^2 - 5 \qquad \frac{4}{x} \qquad x^3 - 2x + 1 \qquad
      \sqrt{x} + 6 \qquad 7 \qquad 2^x
    \]
    \begin{enumerate}[itemsep=2pt]
        \item Which three did you circle, and what rule did you use?
        \item Does the lone $7$ belong with them or not? Defend your answer.
    \end{enumerate}
    \textit{Discussion prompt:} students almost always sort correctly on instinct
    and then struggle to say why. Collect three or four attempted rules on the
    board without judging them --- the lesson's job is to sharpen the instinct into
    a definition. The $7$ is the productive argument: it is a polynomial (degree 0).
\end{skillbox}

\begin{skillbox}[Lesson]{sky}
\begin{multicols}{2}

\textbf{Part 1: What a Polynomial Is} (10 min)

\begin{itemize}
  \item Start informally: add up terms, and every term is a number times a
        whole-number power of $x$.
  \item Take $p(x) = 4x^3 - 7x^2 + 2x - 9$ apart term by term. Point out that
        $2x$ is $2x^1$ and $-9$ is $-9x^0$ --- every term really does fit.
  \item Then write the formal statement:
        $p(x) = a_n x^n + \cdots + a_1x + a_0$, with $n$ a positive integer,
        each $a_i$ real, and $a_n \ne 0$. Read it aloud in English.
  \item Close the loop on the hook: a constant is a polynomial of degree 0.
\end{itemize}

\textbf{Part 2: The Parts, by Name} (10 min)

\begin{itemize}
  \item Degree, leading term, leading coefficient, constant term --- point at each
        one on the projected $p(x)$ and have students point at their copy.
  \item Standard form is descending powers. Rewrite $5x - 2 + 3x^4 - x^2$ together.
  \item Spend a full minute on $-x^3$: the leading coefficient is $-1$, not $1$.
        This single error costs points all unit.
\end{itemize}

\columnbreak

\textbf{Part 3: Spotting the Impostors} (15 min)

\begin{itemize}
  \item This is the highest-value block of the lesson --- do not rush it.
  \item Work the notes table one row at a time, rewriting each impostor to expose
        the problem: $\frac{4}{x} = 4x^{-1}$, $\sqrt{x} = x^{1/2}$.
  \item Land the three-question test: whole-number exponents? $x$ out of every
        denominator? $x$ out of every exponent?
  \item Then defuse the two traps: $\frac{x^2}{5}$ \textit{is} a polynomial
        (dividing by a number is multiplying by $\frac{1}{5}$), and $\sqrt{7}\,x^2$
        \textit{is} a polynomial (the radical is a coefficient, not an exponent).
\end{itemize}

\textbf{Part 4: The Family Picture} (10 min)

\begin{itemize}
  \item Name degrees 0--4. Degrees 0, 1, 2 are old friends; 3 and 4 are new.
  \item Project the three pre-drawn graphs. Smooth (no corners), unbroken (no
        holes, no jumps). Count turns; observe at most $n-1$.
  \item Keep this qualitative --- extrema and end behavior are 2.2 and 2.5. Naming
        a turn on a picture is enough today.
  \item Close with the unit map (notes section 6): what 2.1 through 2.7 will do.
\end{itemize}
\end{multicols}
\end{skillbox}

\begin{skillbox}[Active Monitoring]{redbox}
While students work on guided notes and practice, circulate and check:
\begin{itemize}
  \item \textbf{Common error:} calling $\frac{x^2}{5}$ a non-polynomial because it
        ``has a fraction.'' Ask: ``Where is the $x$ --- top or bottom?''
  \item \textbf{Common error:} leading coefficient of $-x^3$ recorded as $1$.
        Cold-call: ``Read me the whole term, sign included.''
  \item \textbf{Common error:} reading degree off the \textit{first} term written
        rather than the highest power. Insist on standard form first.
  \item \textbf{Common error:} constant term of $-6x^4 + 3x^2 + x$ left blank
        instead of $0$. Ask: ``What is $p(0)$?''
  \item \textbf{Cold-call:} ``Give me an expression that is one small change away
        from being a polynomial. Now make the change.''
\end{itemize}
\end{skillbox}

\begin{skillbox}[Group Work \& Differentiation]{redbox}
Students work on the \textbf{Group Activity: Sorting the Polynomial Family} in
leveled tiers. All three tiers work from the same eight-expression bank (A--H), so
groups can compare answers across tiers.

\begin{itemize}
  \item \textbf{Tier R (Remediate):} Sort all eight as polynomial or not, and give
        the degree of each one that is. Three follow-up counting questions.
  \item \textbf{Tier A (Apply):} Rewrite three out-of-order polynomials in standard
        form and read off degree, leading coefficient, and constant term --- one
        of them with a hidden constant term of $0$.
  \item \textbf{Tier E (Extend):} Name the specific rule each impostor breaks,
        build a polynomial to stated specifications (degree 3, leading coefficient
        $-2$, constant term 5), settle the $\frac{x^3 - 2x}{4}$ argument, and
        explain from the definition why degree $\frac{1}{2}$ is impossible.
\end{itemize}
\end{skillbox}

\begin{skillbox}[Individual Work \& Assessment]{redbox}
\textbf{Exit Ticket} (last 5--7 minutes, individual, collected):
\begin{enumerate}
  \item Give the degree, leading coefficient, and constant term of
        $p(x) = -4x^3 + 7x^2 - x + 9$, and name its degree family.
  \item Decide which of $6x^4 - 2x + 1$ and $3x^2 + \frac{5}{x}$ is not a
        polynomial, with one sentence of reasoning.
  \item Write $5 - 2x + 4x^2$ in standard form.
\end{enumerate}
\textbf{Assessment note:} Item 2's sentence is the one to read closely --- look for
a reason naming the \textit{exponent} or the \textit{denominator}, not just
``it has a fraction.'' A vague reason predicts trouble in 2.6 and in Unit 3, where
the same expression will be a rational function. Item 1's leading coefficient
($-4$, not $4$) is the sign check.
\end{skillbox}

\begin{skillbox}[Reinforcement \& Extension]{goldbox}
\textbf{Homework:} 5 problems --- naming the parts of two polynomials, four
polynomial-or-not judgments with reasons, three standard-form rewrites with degree
families, turning-point counts read off two pre-drawn graphs, and one SOL-style
multiple-choice item.

\textbf{Extension (optional):} A smooth, unbroken graph has exactly 4 turning
points. What is the smallest degree possible, and could the degree be 7 instead?
Students reason with ``at most $n-1$'' in both directions.

\textbf{Preview:} Next lesson (2.1, Quadratic Functions Revisited) zooms in on the
degree-2 members of the family --- vertex, axis of symmetry, and the forms a
quadratic can be written in --- as the model for everything the unit does with
higher degrees.
\end{skillbox}

\begin{teachernote}[Warm-Up]
Five minutes, silent start, no calculators. All four items are Algebra 1--2
maintenance, chosen because today's content sits directly on top of them: exponent
evaluation, like terms, function notation, graph reading. A student who cannot
combine like terms in item 2 will not be able to write a polynomial in standard
form this period --- flag them before the activity, not after.
\end{teachernote}

\begin{teachernote}[Guided Notes]
Sections 1 and 2 move quickly; section 3 is the lesson and deserves the time. Work
the impostor table row by row rather than assigning it, and rewrite each impostor
in exponent form on the board so the disqualifier is visible rather than asserted.
The two ``yes'' traps at the bottom of that table ($\frac{x^2}{5}$ and
$\sqrt{7}\,x^2 - \frac{1}{2}x$) are the rows that produce the most argument --- let
the argument happen. Section 6 is read-only; assign it aloud in the last minute of
notes and do not add to it.
\end{teachernote}

\begin{teachernote}[Group Activity]
Mixed-readiness groups work well since every tier reads the same bank. Tier R
groups should say the disqualifying reason aloud even though the tier does not ask
them to write it. In Tier A, the third row's constant term of $0$ is the item worth
circulating for. If a Tier E pair resolves the Marisol/Rafi argument cleanly, have
them present it in the closing minutes --- it is the same misconception the exit
ticket tests.
\end{teachernote}

\begin{teachernote}[Exit Ticket]
Collected, completion credit only (``mistakes happen, blanks don't''). Sort by item
2's reasoning and re-open the three-question test in the 2.1 warm-up if more than a
handful gave a vague reason. Item 3 takes ten seconds to grade and tells you
whether standard form landed.
\end{teachernote}

\begin{teachernote}[Homework]
Problem 4 is the first item asking students to read a polynomial graph; it wants
counting only, no analysis --- if students start describing where the function
increases, that is fine but it is 2.2's job. Problem 5 is the SOL-style
multiple-choice item. The extension mirrors the ``at most $n-1$'' idea from notes
section 5 and is a fair second try for anyone who missed problem 4.
\end{teachernote}
\end{document}
